\documentclass[11pt]{article}

% Prefix for numedquestion's
\newcommand{\questiontype}{Question}


% Use this if your "written" questions are all under one section
% For example, if the homework handout has Section 5: Written Questions
% and all questions are 5.1, 5.2, 5.3, etc. set this to 5
% Use for 0 no prefix. Redefine as needed per-question.
\newcommand{\writtensection}{0}

\usepackage{amsmath, amsfonts, amsthm, amssymb}  % Some math symbols
\usepackage{mathtools}
\usepackage{algorithm}
\usepackage{algpseudocode}
\usepackage{dsfont}

\newtheorem{theorem}{Theorem}[section]
\theoremstyle{definition}
\newtheorem{definition}{Definition}[section]
\newtheorem{lemma}[theorem]{Lemma}

\usepackage{centernot}
\usepackage{mathtools}

\usepackage{enumitem}

\setlength{\parindent}{0pt}

\begin{document}
\section{Tree Evaluation}
Consider the function:
$$f: \{0,1\}^k \times \{0,1\}^k \rightarrow \{0,1\}^k$$
We aim to find the amount of space needed, given $x_1,...,x_n \in \{0,1\}^k$ for $n=2^k$, the root value of the tree defined where each inner node is calculated as:
$$u.val = f(u.left.val, u.right.val)$$
Consider a tree that has $2^k$ leaf nodes ranging from $x_1,...,x_{2^k}$, the trivial algorithm to evaluate the root node's value would take $O(k \cdot k)$ space (the height of the tree). The trivial algorithm is as follows:
\begin{verbatim}
Evaluate(node)
    If node == leaf:
        return node.value
    Else:
        return f(Evaluate(node.left), Evaluate(node.right))
\end{verbatim}
The space comes from needing to store the value of $f(u)$ (k bits) for each level which we refer to as $h$ for height where $(h = k)$.
\bigskip

To help understand the problem, we imagine that ever node stores the height at which it is in the tree, and for each level of the tree, we have a register $R_i$ where $1 \leq i \leq k$. Each output will be stored in $R_d$, where $d = height(u)$. Now we have:
\begin{verbatim}
    If u is a leaf:
        R1 = u.val
    TEVAL(u.left)
    Rd <- Rd-1
    TEVAL(u.right)
    Rd <- f(Rd-1,Rd)
\end{verbatim}
This eliminates the need for recursion. The space here is $\Theta(h\cdot k)$. We assume that $h = k = \log n$
\section{Cook-Mertz}
A recent result has shown that TEVAL can be solved in $O((k+h)\log k)$ space, or $O(\log n \log \log n)$. This has been applied to show a theorem in 2025 by Ryan Williams that:
$$TIME(t(n)) \subseteq SPACE(\sqrt{t(n)}\cdot polylog(t(n)))$$
\subsection{Formula Evaluation With Minimal Extra Registers}
Imagine we have a formula as follows:
$$F = (2+3) * (6 + 7)$$
To solve this recursively, we can evaluate the formual as a tree. Above, we just reduce the problem to a tree, with leaves being numerica values and the inner nodes representing the operations. We do the same thing where each output is stored in $R_d$ where $d = $ height of $u$:
\begin{verbatim}
                FEVAL(u):
                    If u is leaf:
                        R1 = u.val
                        return
                    FEVAL(u.left)
                    Rd <- Rd-1
                    FEVAL(u.right)
                    If u.op == +:
                        Rd <- Rd + Rd-1
                    Else:
                        Rd <- Rd * Rd-1
\end{verbatim}
The space complexity here will be the depth of the tree, $h$ registers for each level of the tree.
\bigskip

In 1989, Barrington showed that regardless of the height of the tree, this formula can be evaluated with a constant number of registers. Ben-Or and Cleave demonstrated it can be done in 3 registers, which is essentially optimal for the root nodes left, right, and own value. It is done with something called a \textbf{Straight Line Program} that is executed sequentially without any loops or conditional statements. It also uses \textbf{Catalytic Computation}, where memory in recursive calls is somewhat "restored" after the recursive calls are completed. It uses the following definition for a new operation called offset.
$$offset(u,R_i,R_j,R_k, "sign")$$
Where sign is $1/-1$. The Goal for each Offset call is to return the following without changing $R_i$ or $R_k$:
$$R_j \leftarrow R_j+(sign)R_i \cdot FEVAL_u(x)$$
We have the following for the case where $sign = 1$.
\bigskip

If u == leaf: $R_j \leftarrow R_j + (sign)R_i(u.value)$
\bigskip

If u.op == +: $R_j \leftarrow R_j + R_i(FEVAL_{u.left}(x) + FEVAL_{u.right}(x))$ We compute this by making the following calls:
    \begin{itemize}
        \item $offset(u.left,R_i,R_j,R_k,Sign)$ Which causes the registers value to update to $R_j \leftarrow R_j + R_i \cdot (FEVAL_{u.left}(x))$
        \item $offset(u.right,R_i,R_j,R_k,Sign)$ Further updates it to \\$R_j \leftarrow R_j + R_i \cdot (FEVAL_{u.left}(x)) + R_i \cdot (FEVAL_{u.right}(x))$
    \end{itemize}
    For u.op == $*$, we want the offset call for $u$ to update registers so that $$R_j \leftarrow R_j + R_i \cdot FEVAL_{u.left}(x) * FEVAL_{u.right}(x)$$
    First, we will rename $FEVAL_{u.left}(x)$ as $f$ and $FEVAL_{u.right}(x)$ as $g$ for simplicity. Our goal is to go from:
    $$(R_i,R_j,R_k) \rightsquigarrow (R_i, R_j + R_i \cdot f \cdot g, R_k)$$
    We consider the following operations that will help achieve this:
    \[
    \begin{array}{|c|c|}
    \hline
    \text{Operation} & \text{Updated registers} \\
    \hline
    R_j \leftarrow R_j - R_k \cdot g & (R_i,\, R_j - R_k \cdot g,\, R_k) \\
    \hline
    R_k \leftarrow R_k + R_i \cdot f & (R_i, R_j - R_k \cdot g, R_k + R_i \cdot f)\\
    \hline
    R_j \leftarrow R_j + R_k \cdot g & \begin{aligned}
        &(R_i,(R_j - R_k\cdot g) + (R_k + R_i \cdot f)\cdot g, R_k + R_i \cdot f)\\
        &= (R_i, R_j + R_i\cdot f\cdot g, R_k + R_i\cdot f)
    \end{aligned}\\
    \hline
    R_k \leftarrow R_k - R_i \cdot f & (R_i, R_j + R_i \cdot f \cdot g, R_k)\\
    \hline
    \end{array}
    \]

\end{document}